\documentclass[11pt,a4paper]{article}
\usepackage[margin=2.5cm]{geometry}
\usepackage{fontspec}
\usepackage[brazilian]{babel}
\usepackage{hyperref}
\setlength{\parindent}{0pt}
\setlength{\parskip}{0.6em}
\setlength{\emergencystretch}{2em}
\begin{document}
\begin{center}
{\large Universidade de Passo Fundo}\\
Curso de Ciência da Computação\\[1em]
{\Large Declaração de uso de inteligência artificial}\\
\textbf{Minuta para revisão e assinatura -- não assinada}
\end{center}

\textbf{Trabalho:} Sistema de Monitoramento Agrícola por Imagens de Satélite.\\
\textbf{Autor:} Eduardo Steffens Hoppen.\\
\textbf{Orientador:} Prof. Rafael Rieder.

Esta minuta segue os quatro campos do anexo da IN nº 03/2026/PROACAD, de 28 de outubro de 2026. Seu preenchimento não certifica autorização institucional, aprovação ética ou validação final da revisão atual.

Declaro, para os devidos fins, que este trabalho utilizou ferramentas de Inteligência Artificial nas seguintes condições:

\section*{1. Ferramentas utilizadas}
OpenCode, com identificadores locais \nolinkurl{openai/gpt-6.1-terra} e \nolinkurl{openai/gpt-6.1-sol}. A instalação consultada informa OpenCode v2.3.0; não foram estabelecidas as versões de cada sessão histórica. Declara-se também o uso denominado ``ChatGPT 6 -- Luna'', sem identificação documental de versão ou vínculo com as sessões recuperadas. Não se presume correspondência entre esses identificadores e nomes comerciais oficiais.

\section*{2. Finalidade do uso}
Auxílio na elaboração de especificações e tarefas; geração e revisão de código e testes; documentação técnica; execução de verificações por ferramentas; pesquisa, organização e conferência assistida de referências; revisão científica e matemática; geração e reformulação substancial de trechos do artigo; revisão gramatical e didática; tradução do resumo; elaboração de diagramas e scripts; análise quantitativa assistida, codificação e agrupamento qualitativo; compilação e inspeção visual assistida.

A IA não substitui o contato direto do pesquisador com as fontes nem sua responsabilidade sobre as decisões de arquitetura, o conteúdo e as conclusões. Os resultados técnicos gerados pelas ferramentas e as aprovações humanas são registrados separadamente.

\section*{3. Etapas da atividade acadêmica}
Planejamento e especificação de funcionalidades; implementação e revisão do artefato; preparação e execução de testes; documentação de demonstrações manuais; tratamento e análise do arquivo do questionário; elaboração, revisão bibliográfica, científica e editorial do artigo; preparação de figuras, arquivos LaTeX e PDF; revisão final ainda sujeita à aprovação do autor.

As decisões técnicas e arquiteturais e a supervisão permaneceram sob responsabilidade do pesquisador, que conferiu as treze categorias temáticas e as interpretações qualitativas. A leitura e aprovação das reformulações científicas e editoriais desta versão permanecem necessárias antes da assinatura.

\section*{4. Prompts relevantes}
Os seguintes excertos são literais de mensagens recuperadas do projeto no OpenCode, e não reconstruções de prompts:

\begin{quote}
``Antes de consolidar os resultados, verifique a procedência dos registros e eventuais correspondências com bases sintéticas existentes no projeto. Separe registros comprovadamente reais, sintéticos e de procedência não confirmada.''
\end{quote}
Mensagem \texttt{msg\_1e0f1bee5001QnNu3YHb8uvtC7}, sobre a avaliação de usuários.

\begin{quote}
``Não atribuir esses limiares a de Jong et al.''
\end{quote}
Mensagem de execução final, seção de regressão linear, \texttt{msg\_1e1f9356f001lL77wG8k7hQ4Rx}.

A terceira rodada de consolidação também foi orientada pela instrução literal: ``O artigo deve comunicar a pesquisa realizada, e não narrar o processo de revisão científica pelo qual passou.'' A mensagem integral desta solicitação permanece na conversa do projeto.

O registro complementar \texttt{REGISTRO\_USO\_IA.md} relaciona as sessões, as especificações, as finalidades e as evidências. Os exemplos acima não representam um histórico completo. Os prompts integrais devem ser disponibilizados ao orientador por canal adequado, após revisão de privacidade; não se publicam automaticamente dados inéditos ou respostas individuais.

\section*{Declaração autoral para assinatura}
O texto abaixo somente deve ser subscrito após a leitura e validação efetivas da versão entregue:

Declaro que todo o conteúdo gerado com auxílio de IA foi criticamente revisado, validado e, quando necessário, substancialmente reformulado por mim, e que assumo integral responsabilidade pela originalidade, veracidade e integridade científica do trabalho apresentado.

Declaro, ainda, ter conhecimento de que a omissão ou a falsificação desta declaração configura infração e poderá ser equiparada à má conduta científica, nos termos do Código de Ética, do Regimento Geral da Universidade de Passo Fundo e da Portaria CNPq nº 2.664/2026.

\vspace{1em}
Passo Fundo, \rule{1cm}{0.4pt} de \rule{4cm}{0.4pt} de \rule{1.5cm}{0.4pt}.

\vspace{2em}
\begin{center}
\rule{10cm}{0.4pt}\\
Eduardo Steffens Hoppen\\
Nome e assinatura
\end{center}
\end{document}
